% =============================================================================
%  lec15.tex —— 第 15 讲 稳定 ∞-范畴（底本 PDF 第 117–123 页）
% =============================================================================
\lecture{15}{稳定 ∞-范畴}
\begin{center}\small\color{black!55}（底本 PDF 第 117–123 页）\end{center}

本讲定义并研究\textbf{稳定 $\infty$-范畴}，它们恰是形如 $\mathrm{Sptr}(\cC)$ 的 $\infty$-范畴。

\section{稳定性}

\begin{definition}\label{L15-1.1}
设 $\cC$ 是带点 $\infty$-范畴。$\cC$ 中的一个\textbf{三角}（triangle）是下图所示的一个图 $\Delta^{1}\times\Delta^{1}\to\cC$：
\begin{equation}\label{L15-1.1eq}
\begin{tikzcd}[column sep=3.6em,row sep=3em]
X \arrow[r,"f"] \arrow[d] & Y \arrow[d,"g"] \\
0 \arrow[r] & Z
\end{tikzcd}
\end{equation}
其中 $0\in\cC$ 是零对象。若这样的三角是拉回方块，则称它为\textbf{纤维序列}；若它是推出方块，则称它为\textbf{余纤维序列}。

对态射 $f:X\to Y$，$f$ 的\textbf{余纤维}是一个纤维序列 \eqref{L15-1.1eq}。对偶地，对态射 $g:Y\to Z$，$g$ 的\textbf{纤维}是一个余纤维序列 \eqref{L15-1.1eq}。
\end{definition}

\yizhu{原文此处把“纤维序列”与“余纤维序列”对调了：按定义，$f$ 的余纤维应配余纤维序列 \eqref{L15-1.1eq}（即推出方块），$g$ 的纤维应配纤维序列 \eqref{L15-1.1eq}（即拉回方块）。译文照录原文，以便与底本对勘。\textbf{[译者按（无外部依据）]}}

\nitem\label{L15-1.2} 我们常滥用记号，把三角写作 $X\xrightarrow{f}Y\xrightarrow{g}Z$。

\begin{warning}\label{L15-1.3}
三角 \eqref{L15-1.1eq} 的数据并不由链 $X\xrightarrow{f}Y\xrightarrow{g}Z$ 决定，即使在同伦意义下也不由它决定。事实上，知道这样一个三角，等价于知道 $g\circ f$ 的一个零伦；而后者即使在同伦意义下也不唯一。
\end{warning}

\nitem\label{L15-1.4} 然而注意，纤维序列 \eqref{L15-1.1eq} 由态射 $g$ 本质唯一地决定；对偶地，余纤维序列 \eqref{L15-1.1eq} 由态射 $f$ 本质唯一地决定。因此，只要把 \eqref{L15-1.1eq} 作为其定义的一部分数据，我们就可以使用记号
\[
\operatorname{Fib}(g),\qquad \operatorname{Cofib}(f)\in\cC .
\]

\begin{definition}\label{L15-1.5}
称 $\infty$-范畴 $\cC$ 是\textbf{稳定的}，若它满足以下条件：
\begin{itemize}[leftmargin=2.2em,itemsep=0.1em]
\item 它是带点的；
\item $\cC$ 中任一态射都有纤维与余纤维；
\item $\cC$ 中的一个三角是纤维序列，当且仅当它是余纤维序列。
\end{itemize}
\end{definition}

\nitem\label{L15-1.6} 对稳定 $\infty$-范畴，我们可以使用\textbf{纤维—余纤维序列}（fiber-cofiber sequence）这一说法。

\begin{exercise}\label{L15-1.7}
找出所有这样的普通范畴：把它们视为 $\infty$-范畴时是稳定的。
\end{exercise}

\begin{proposition}\label{L15-1.8}
设 $\cC$ 是带点 $\infty$-范畴，且既有有限极限又有有限余极限。则以下条件等价。
\begin{enumerate}[label=(\roman*),leftmargin=2.4em,itemsep=0.1em]
\item 函子 $\Sigma:\cC\to\cC$ 全忠实。
\item $\cC$ 中任一形如 $X\to 0\to Z$ 的余纤维序列也是纤维序列。
\item $\cC$ 中任一余纤维序列都是纤维序列。
\item $\cC$ 中任一推出方块都是拉回方块。
\end{enumerate}
\end{proposition}

\begin{corollary}\label{L15-1.9}
设 $\cC$ 是带点 $\infty$-范畴，且既有有限极限又有有限余极限。则以下条件等价。
\begin{enumerate}[label=(\roman*),leftmargin=2.4em,itemsep=0.1em]
\item 函子 $\Sigma:\cC\rightleftarrows\cC:\Omega$ 都是等价。
\item $\cC$ 中形如 $X\to 0\to Z$ 的三角是余纤维序列，当且仅当它是纤维序列。
\item $\cC$ 中的一个三角是余纤维序列，当且仅当它是纤维序列。
\item $\cC$ 中的一个方块是推出，当且仅当它是拉回。
\end{enumerate}
\end{corollary}

\begin{proof}[命题~\ref{L15-1.8} 的证明]
蕴涵 (ii)$\Leftarrow$(iii)$\Leftarrow$(iv) 是明显的。剩下要证 (i)$\Rightarrow$(iv)。设 $\Sigma$ 全忠实。由于 $\Omega$ 是 $\Sigma$ 的右伴随，我们有 $\id_{\cC}\simeq\Omega\circ\Sigma$。对推出方块
\[
\begin{tikzcd}[column sep=3.8em,row sep=2.8em]
W \arrow[r] \arrow[d] & X \arrow[d] \\
Y \arrow[r] & Z,
\end{tikzcd}
\]
我们需要证明 $W\to X\times_{Z}Y$ 可逆。考虑下述交换图
\[
\begin{tikzcd}[column sep=3.4em,row sep=2.8em]
W \arrow[r] \arrow[d] & X \arrow[r] \arrow[d] & 0 \arrow[dd] \\
Y \arrow[r] \arrow[d] & Z \arrow[dr, dotted] & \\
0 \arrow[rr] & & \Sigma W
\end{tikzcd}
\]
这里内层与外层方块都是推出方块，故存在本质唯一的虚线箭头 $Z\to\Sigma W$ 使上图交换。由拉回的函子性，我们得到态射 $X\times_{Z}Y\to 0\times_{\Sigma W}0\simeq\Omega\Sigma W$，它填入下述交换图
\[
\begin{tikzcd}[column sep=3.8em,row sep=3em]
W \arrow[r] \arrow[d,"\simeq"] & X\times_{Z}Y \arrow[d,"\simeq"] \arrow[dl, dotted] \\
\Omega\Sigma W \arrow[r] & \Omega\Sigma X\times_{\Omega\Sigma Z}\Omega\Sigma Y,
\end{tikzcd}
\]
其中竖直态射是同构，因为 $\id_{\cC}\simeq\Omega\circ\Sigma$。由同构的六选二性质，我们得到所需的 $X\xrightarrow{\ \simeq\ }X\times_{Z}Y$。
\end{proof}

\begin{proposition}\label{L15-1.10}
设 $\cC$ 是带点 $\infty$-范畴。以下条件等价。
\begin{enumerate}[label=(\alph*),leftmargin=2.4em,itemsep=0.1em]
\item $\infty$-范畴 $\cC$ 是稳定的。
\item $\infty$-范畴 $\cC^{\op}$ 是稳定的。
\item $\infty$-范畴 $\cC$ 有有限余极限，且 $\Sigma:\cC\to\cC$ 是等价。
\item $\infty$-范畴 $\cC$ 有有限极限，且 $\Omega:\cC\to\cC$ 是等价。
\item $\infty$-范畴 $\cC$ 有有限余极限与有限极限，且 $\cC$ 中的一个方块是推出方块，当且仅当它是拉回方块。
\item $\infty$-范畴 $\cC$ 有有限极限，且 $\Omega^{\infty}:\mathrm{Sptr}(\cC)\to\cC$ 是等价。
\end{enumerate}
\end{proposition}

\begin{proof}
等价 (a)$\Leftrightarrow$(b) 是明显的。等价 (d)$\Leftrightarrow$(f) 上一讲已证。

只需再证 (c)$\Leftrightarrow$(a)$\Leftrightarrow$(e)，因为 (d)$\Leftrightarrow$(b) 可由过渡到反范畴得到。

设 $\cC$ 有有限余极限与有限极限，则由命题~\ref{L15-1.8} 中的 (i)$\Leftrightarrow$(iii)$\Leftrightarrow$(iv)（及其对偶版本）可得 (c)$\Leftrightarrow$(a)$\Leftrightarrow$(e)。于是只需再证 (a) 或 (c) 蕴涵 $\cC$ 有有限余极限与有限极限。

对 (a)，我们只需证明稳定 $\infty$-范畴有余等化子。这由下文练习~\ref{L15-2.13} 得出。

对 (c)，我们可以用 $\iota:\cC\to\Ind(\cC)$ 把 $\cC$ 嵌入一个可表现 $\infty$-范畴。注意 $\Ind(\cC)$ 是带点的，因为 $\iota$ 既保持又检测有限余极限与有限极限。此外，$\Sigma_{\Ind(\cC)}$ 可等同于 $\Ind(\Sigma_{\cC}):\Ind(\cC)\to\Ind(\cC)$，因而也是等价。由前述讨论，$\Ind(\cC)$ 满足命题中的全部性质。特别地，$\Ind(\cC)$ 中任一推出方块都是拉回方块。由于 $\iota$ 既保持又检测有限余极限与有限极限，同样的结论对 $\cC$ 成立。特别地，$\cC$ 有拉回，从而有所需的全部有限极限。
\end{proof}

\begin{corollary}\label{L15-1.11}
设 $\cC$ 是有有限极限的带点 $\infty$-范畴。则 $\infty$-范畴 $\mathrm{Sptr}(\cC)$ 是稳定的。
\end{corollary}

\begin{exercise}\label{L15-1.12}
设 $\cC$ 是稳定 $\infty$-范畴。则 $f:X\to Y$ 是同构，当且仅当 $\operatorname{Fib}(f)$ 是零对象，当且仅当 $\operatorname{Cofib}(f)$ 是零对象。
\end{exercise}

\begin{exercise}\label{L15-1.13}
设 $f:X\to Y$ 是稳定 $\infty$-范畴 $\cC$ 中的态射。证明 $\Sigma\operatorname{Fib}(f)\simeq\operatorname{Cofib}(f)$。
\end{exercise}

\section{稳定 $\infty$-范畴的同伦范畴}

\nitem\label{L15-2.1} 本节中，设 $\cC$ 是稳定 $\infty$-范畴。

\begin{exercise}\label{L15-2.2}
证明典范态射 $X\sqcup Y\to X\times Y$ 可逆。提示：
\begin{center}
\begin{tikzcd}[column sep=4em,row sep=2.8em]
X\sqcup Y \arrow[r] \arrow[d] & X\sqcup 0 \arrow[d] \\
0\sqcup Y \arrow[r] & 0\sqcup 0
\end{tikzcd}
\end{center}
是推出方块。
\end{exercise}

\nitem\label{L15-2.3} 由于 $X\sqcup Y$ 与 $X\times Y$ 之间有典范等价，我们用 $X\oplus Y$ 同时表示二者。

\begin{exercise}\label{L15-2.4}
设 $f,g:X\rightrightarrows Y$ 是两个态射。证明复合
\[
X\to X\oplus X\xrightarrow{(f,g)}Y\oplus Y\to Y
\]
给出 $\pi_{0}\operatorname{Maps}_{\cC}(X,Y)$ 上一个良定义的二元运算。我们把上述复合记为 $f+g$。
\end{exercise}

\begin{exercise}\label{L15-2.5}
设 $f,g:X\rightrightarrows Y$ 是两个态射。证明上述二元运算与阿贝尔群
\[
\pi_{0}\operatorname{Maps}_{\cC}(X,Y)\simeq\pi_{0}\operatorname{Maps}_{\cC}(X,\Omega^{2}\Sigma^{2}Y)\simeq\pi_{0}\Omega^{2}\operatorname{Maps}_{\cC}(X,\Sigma^{2}Y)\simeq\pi_{2}\operatorname{Maps}_{\cC}(X,\Sigma^{2}Y)
\]
上的加法运算一致。
\end{exercise}

\begin{exercise}\label{L15-2.6}
设 $\sigma:\Delta^{1}\times\Delta^{1}\to\cC$ 是形如
\begin{center}
\begin{tikzcd}[column sep=3.6em,row sep=2.8em]
X \arrow[r] \arrow[d] & 0 \arrow[d] \\
0 \arrow[r] & \Sigma Y
\end{tikzcd}
\end{center}
的图，$\sigma'$ 是它的转置。由推出的泛性质，$\sigma$ 诱导一个态射 $\eta_{\sigma}:X\to\Omega\Sigma Y\simeq Y$，它是 $\pi_{0}\Hom_{\cC}(X,Y)$ 中一个良定义的元素。证明 $\eta_{\sigma}+\eta_{\sigma'}=0$。
\end{exercise}

\begin{corollary}\label{L15-2.7}
同伦范畴 $\mathrm{h}\cC$ 是加法范畴。
\end{corollary}

\nitem\label{L15-2.8} 此后我们记 $X[n]\defeq\Sigma^{n}X$，其中当 $n<0$ 时取 $\Sigma^{n}\defeq\Omega^{-2}$。注意这些对象在同伦意义下良定义。

\yizhu{原文此处作 $\Sigma^{n}\defeq\Omega^{-2}$；按记号，当 $n<0$ 时应为 $\Sigma^{n}\defeq\Omega^{-n}$。译文照录原文，以便与底本对勘。\textbf{[译者按（无外部依据）]}}

\begin{definition}\label{L15-2.9}
设 $X\xrightarrow{f}Y\xrightarrow{g}Z\xrightarrow{h}X[1]$ 是 $\mathrm{h}\cC$ 中的一条链。称它是 $\mathrm{h}\cC$ 中的一个\textbf{特异三角}（\mbox{distinguished} triangle），若存在一个图 $\Delta^{1}\times\Delta^{2}\to\cC$ 如下图所示
\[
\begin{tikzcd}[column sep=3.6em,row sep=2.8em]
X \arrow[r,"\tilde{f}"] \arrow[d] & Y \arrow[r] \arrow[d,"\tilde{g}"] & 0 \arrow[d] \\
0' \arrow[r] & Z \arrow[r] & W
\end{tikzcd}
\]
使得
\begin{itemize}[leftmargin=2.2em,itemsep=0.1em]
\item 对象 $0$ 与 $0'$ 是零对象；
\item 态射 $\tilde{f}$ 与 $\tilde{g}$ 分别提升 $f$ 与 $g$；
\item 外方块是纤维—余纤维序列；
\item 复合 $Z\to W\xrightarrow{\ \simeq\ }X[1]$（它在同伦意义下良定义）提升 $h$。
\end{itemize}
\end{definition}

\yizhu{“distinguished triangle”在中文文献中或直译作“\mbox{distinguished} 三角”，或意译作“特异三角”“正合三角”。译本取“特异三角”，以与“三角范畴”“三角函子”的译名配套：它指的是 $\mathrm{h}\cC$ 中被选出的那一类三角，而非任意的三角。\textbf{[来源：B]}}

\begin{theorem}[HA.1.1.2.14]\label{L15-2.10}
上述平移函子与特异三角的选取，使 $\mathrm{h}\cC$ 成为一个三角范畴。
\end{theorem}

\begin{samepage}
\begin{exercise}\label{L15-2.11}
若改用形如
\begin{center}
\begin{tikzcd}[column sep=3.6em,row sep=2.4em]
X \arrow[r] \arrow[d] & 0 \arrow[d] \\
Y \arrow[r] \arrow[d] & Z \arrow[d] \\
0' \arrow[r] & W
\end{tikzcd}
\end{center}
的图来定义 $\mathrm{h}\cC$ 中的特异三角，会发生什么？
\end{exercise}
\end{samepage}

\begin{exercise}\label{L15-2.12}
对 $X\in\cC$，构造一个纤维—余纤维序列
\[
X\xrightarrow{(\id,-\id)}X\oplus X\xrightarrow{(\id,\id)}X.
\]
\end{exercise}

\begin{exercise}\label{L15-2.13}
证明 $f,g:X\rightrightarrows Y$ 的余等化子典范等价于 $\operatorname{Cofib}(f-g)$。提示：
\begin{center}
\begin{tikzcd}[column sep=3.6em,row sep=2.8em]
X \arrow[r,"{(\id,-\id)}"] \arrow[d] & X\oplus X \arrow[r,"{(f,g)}"] \arrow[d,"{(\id,\id)}"] & Y \arrow[d] \\
0 \arrow[r] & X \arrow[r] & \bullet
\end{tikzcd}
\end{center}
\end{exercise}

\section{映射谱}

\begin{construction}\label{L15-3.1}
设 $\cC$ 是稳定 $\infty$-范畴。对 $X,Y\in\cC$，定义
\[
\underline{\Hom}(X,Y)_{n}\defeq\operatorname{Maps}(X,Y[n]).
\]
注意对任意 $n\ge 0$，有同构
\[
\alpha_{n}:\operatorname{Maps}(X,Y[n])\simeq\operatorname{Maps}(X,\Omega Y[n+1])\simeq\Omega\operatorname{Maps}(X,Y[n+1]).
\]
设 $\underline{\Hom}(X,Y)\in\mathrm{Sptr}$ 是由空间 $\{\underline{\Hom}(X,Y)_{n}\}$ 与同构 $\alpha_{n}$ 给出的谱。我们称它为 $X$ 与 $Y$ 之间的\textbf{映射谱}（mapping spectrum）。
\end{construction}

\begin{remark}\label{L15-3.2}
上述谱 $\underline{\Hom}(X,Y)$ 在同伦意义下良定义。在后续各讲中，我们将给 $\mathrm{Sptr}$ 配上对称幺半结构，使得任何稳定 $\infty$-范畴 $\cC$ 都典范地以 $\mathrm{Sptr}$ 充实。
\end{remark}

\begin{definition}\label{L15-3.3}
设 $\cC$ 是稳定 $\infty$-范畴。对 $X,Y\in\cC$，定义
\[
\operatorname{Ext}^{n}(X,Y)\defeq\pi_{0}\underline{\Hom}(X,Y)_{n}\simeq\pi_{0}\operatorname{Maps}(X,Y[n]),
\]
并称之为 $X$ 与 $Y$ 之间的\textbf{第 $n$ 个扩张群}（extension group）。
\end{definition}

\nitem\label{L15-3.4} 注意扩张群只依赖 $X$ 与 $Y$ 在三角范畴 $\mathrm{h}\cC$ 中的像。众所周知，对 $\mathrm{h}\cC$ 中任一特异三角 $Y_{0}\to Y_{1}\to Y_{2}\to Y_{0}[1]$，我们有长正合序列
\[
\cdots\to\operatorname{Ext}^{n}(X,Y_{0})\to\operatorname{Ext}^{n}(X,Y_{1})\to\operatorname{Ext}^{n}(X,Y_{2})\to\operatorname{Ext}^{n+1}(X,Y_{0})\to\cdots.
\]

\begin{exercise}\label{L15-3.5}
设 $\cC$ 是有小余极限的稳定 $\infty$-范畴，$X\in\cC$ 是对象。证明以下条件等价：
\begin{itemize}[leftmargin=2.2em,itemsep=0.1em]
\item 对象 $X$ 是紧的，即 $\operatorname{Maps}_{\cC}(X,-):\cC\to\Spc$ 保持小滤过余极限。
\item 函子 $\underline{\Hom}(X,-):\cC\to\mathrm{Sptr}$ 保持小余极限。
\item 函子 $\underline{\Hom}(X,-):\cC\to\mathrm{Sptr}$ 保持小余积。
\end{itemize}
\end{exercise}

\begin{warning}\label{L15-3.6}
上文中函子 $\operatorname{Maps}_{\cC}(X,-)$ 不保持一般的有限余极限，例如悬挂。
\end{warning}

\section{正合函子}

\nitem\label{L15-4.1} 命题~\ref{L15-1.10} 与练习~\ref{L15-2.13} 蕴涵下述结果。

\begin{theoremdef}\label{L15-4.2}
设 $F:\cC\to\cC'$ 是稳定 $\infty$-范畴之间的函子。以下条件等价。
\begin{itemize}[leftmargin=2.2em,itemsep=0.1em]
\item 函子 $F$ 保持零对象与纤维—余纤维序列。
\item 函子 $F$ 是左正合的，即保持有限极限。
\item 函子 $F$ 是右正合的，即保持有限余极限。
\end{itemize}
若 $F$ 满足上述条件，则称 $F$ 是\textbf{正合的}。
\end{theoremdef}

\yizhu{原文作 Proposition-Definition（命题—定义），即把命题与定义合成一条陈述；本译本只提供合并的“定理—定义”环境，故排作“定理—定义 4.2”，其类别仍按原文理解。\textbf{[译者按（无外部依据）]}}

\begin{exercise}\label{L15-4.3}
设 $F:\cC\to\cC'$ 是稳定 $\infty$-范畴之间的正合函子。证明 $\mathrm{h}F:\mathrm{h}\cC\to\mathrm{h}\cC'$ 自然地带有三角范畴之间正合函子\footnote{又称三角函子。警告：成为三角函子是\textbf{结构}而不是性质。}的结构。
\end{exercise}

\begin{definition}\label{L15-4.4}
设 $\Cat_{\infty}^{\mathrm{ex}}\subseteq\Cat_{\infty}$ 是由小稳定 $\infty$-范畴及其间的正合函子构成的子 $\infty$-范畴。
\end{definition}

\begin{exercise}\label{L15-4.5}
设 $\cC$ 是稳定 $\infty$-范畴。证明 $\Sigma$ 与 $\Omega$ 是正合的。它们诱导的三角函子是什么？
\end{exercise}

\begin{exercise}\label{L15-4.6}
设 $F:\cC\to\cC'$ 是稳定 $\infty$-范畴之间的正合函子。证明 $F$ 是保守的，当且仅当 $F$ 检测零对象。
\end{exercise}

\section{封闭性质}

\begin{exercise}\label{L15-5.1}
设 $\cC$ 是稳定 $\infty$-范畴，$K$ 是单纯集。证明 $\Fun(K,\cC)$ 是稳定的。
\end{exercise}

\begin{exercise}\label{L15-5.2}
设 $\cC$ 与 $\cD$ 是稳定 $\infty$-范畴。证明 $\Fun_{\mathrm{ex}}(\cC,\cD)$ 是稳定的。
\end{exercise}

\begin{exercise}\label{L15-5.3}
设 $\cC$ 与 $\cD$ 是可表现 $\infty$-范畴，且 $\cD$ 是稳定的。证明 $\mathrm{LFun}(\cC,\cD)$ 是稳定的。
\end{exercise}

\begin{exercise}\label{L15-5.4}
设 $\cC$ 是小稳定 $\infty$-范畴，则 $\Ind(\cC)$ 是稳定的。
\end{exercise}

\begin{exercise}\label{L15-5.5}
设 $\cC$ 是稳定 $\infty$-范畴，则 $\cC^{\mathrm{cpt}}$ 是稳定的。特别地，$\mathrm{Sptr}^{\mathrm{fin}}$ 是稳定的。
\end{exercise}

\begin{exercise}\label{L15-5.6}
设 $\cC$ 是稳定 $\infty$-范畴，证明 $\cC$ 的幂等完备化（idempotent completion）也是稳定的。
\end{exercise}

\begin{exercise}\label{L15-5.7}
设 $\cC$ 是稳定 $\infty$-范畴，$\mathrm{PShv}(\cC)$ 是稳定的吗？
\end{exercise}

\begin{theorem}\label{L15-5.8}
$\infty$-范畴 $\Cat_{\infty}^{\mathrm{ex}}$ 有小极限，且包含函子 $\Cat_{\infty}^{\mathrm{ex}}\to\Cat_{\infty}$ 保持并检测小极限。
\end{theorem}

\begin{proof}[证明梗概]
我们只需证明：对任一以 $i\mapsto\cC_{i}$ 给出的小图 $K\to\Cat_{\infty}$，只要每个 $\cC_{i}$ 是稳定的、每个连接函子 $\cC_{i}\to\cC_{j}$ 是正合的，就有
\begin{itemize}[leftmargin=2.2em,itemsep=0.1em]
\item 极限 $\infty$-范畴 $\cC\defeq\lim_{i}\cC_{i}$ 是稳定的；
\item 求值函子 $\cC\to\cC_{i}$ 是正合的。
\end{itemize}
两项断言都可用 $\cC$ 中对象与映射空间的显式描述来验证。
\end{proof}

\nitem\label{L15-5.9} 类似地，可以证明下述结果。

\begin{theorem}\label{L15-5.10}
$\infty$-范畴 $\Cat_{\infty}^{\mathrm{ex}}$ 有小滤过余极限，且包含函子 $\Cat_{\infty}^{\mathrm{ex}}\to\Cat_{\infty}$ 保持并检测小滤过余极限。
\end{theorem}

\section{$\mathrm{Sptr}$ 的泛性质}

\nitem\label{L15-6.1} 下述结果蕴涵：$\mathrm{Sptr}$ 是在小余极限下由单个对象自由生成的稳定 $\infty$-范畴。

\begin{exercise}\label{L15-6.2}
设 $\cD$ 是可表现稳定 $\infty$-范畴。证明在 $\mathcal{S}\in\mathrm{Sptr}$ 处的求值诱导一个等价
\[
\mathrm{LFun}(\mathrm{Sptr},\cD)\xrightarrow{\ \simeq\ }\cD.
\]
提示：证明 $\mathrm{RFun}(\cD,\mathrm{Sptr})\xrightarrow{\ \Omega^{\infty}\circ-\ }\mathrm{RFun}(\cD,\Spc)$ 是等价。
\end{exercise}

\begin{exercise}\label{L15-6.3}
设 $\cD$ 是可表现稳定 $\infty$-范畴。证明在 $\mathcal{S}\in\mathrm{Sptr}^{\mathrm{fin}}$ 处的求值诱导一个等价
\[
\Fun_{\mathrm{ex}}(\mathrm{Sptr}^{\mathrm{fin}},\cD)\xrightarrow{\ \simeq\ }\cD.
\]
\end{exercise}

\appendixsection{无模型的三角范畴}

\nitem\label{L15-A.1} 存在一些三角范畴，它们不是任何稳定 $\infty$-范畴的同伦范畴。

\nitem\label{L15-A.2} 稳定 $\infty$-范畴的同伦范畴之间存在一些正合函子，它们不来自这些稳定 $\infty$-范畴之间的正合函子。

\nitem\label{L15-A.3} \textbf{建议阅读}\quad [MSS07]。

\subsection*{参考文献}
\par\noindent\hangindent=5.2em\hangafter=1 [MSS07]\hspace{0.6em}Fernando Muro, Stefan Schwede, and Neil Strickland. \emph{Triangulated categories without models}. \emph{Inventiones mathematicae}, 170(2):231--241, 2007.
